\subsection{Modes of Probabilistic Convergence}

Pointwise convergence is too strict for most random sequences. The useful replacement is to ask whether probabilities, laws, or moments of the errors become small.

\begin{definition} \label{def:probability-and-distribution-convergence} \label{def:lp-quadratic-and-almost-sure-convergence} \hlt{Modes of Probabilistic Convergence}\\
Let $(X_n)_{n\geq1}$ and $X$ be real-valued random variables with distribution functions $F_n$ and $F$.
\begin{enumerate}[label=(\roman*)]
\setlength{\itemsep}{0pt}
\item \hlt{Convergence in probability}: $X_n\xrightarrow{\mathbb{P}}X$ iff $\mathbb{P}(\abs{X_n-X}>\varepsilon)\to0$ for every $\varepsilon>0$.
\item \hlt{Convergence in distribution}: $X_n\rightsquigarrow X$ iff $F_n(t)\to F(t)$ at every continuity point $t$ of $F$.
\item \hlt{Convergence in $L^p$}: for $p\geq1$, $X_n\xrightarrow{L^p}X$ iff $\mathbb{E}\abs{X_n-X}^p\to0$.
\item \hlt{Convergence in quadratic mean}: $X_n\xrightarrow{qm}X$ iff $X_n\xrightarrow{L^2}X$.
\item \hlt{Almost-sure convergence}: $X_n\xrightarrow{\mathrm{a.s.}}X$ iff $\mathbb{P}(\{\omega\in\Omega:X_n(\omega)\to X(\omega)\})=1$.
\item \hlt{Constant limits}: if $X=c$ almost surely for a constant $c$, then $X_n\xrightarrow{\mathbb{P}}c$, $X_n\rightsquigarrow c$, $X_n\xrightarrow{L^p}c$, $X_n\xrightarrow{qm}c$, and $X_n\xrightarrow{\mathrm{a.s.}}c$ denote convergence to that constant random variable in the corresponding modes.
\end{enumerate}
\end{definition}

\begin{example} \label{ex:shrinking-normal-converges-to-point-mass} \hlt{A Shrinking Normal Law}\\
Let $X_n\sim N(0,1/n)$, and let $F$ be the distribution function of $\delta_0$.\\
Since $\sqrt{n}X_n\sim N(0,1)$, $F_n(t)=\mathbb{P}(X_n\leq t)=\Phi(\sqrt{n}t)$, so $F_n(t)\to0$ for $t<0$ and $F_n(t)\to1$ for $t>0$.\\
Hence $F_n(t)\to F(t)$ at every continuity point $t\neq0$ of $F$, and therefore $X_n\rightsquigarrow0$.
\end{example}

\begin{theorem} \label{thm:probabilistic-convergence-implications} \hlt{Basic Implications Between Convergence Modes}\\
Let $X_n$ and $X$ be real-valued random variables with distribution functions $F_n$ and $F$, and let $c\in\mathbb{R}$.
\begin{enumerate}[label=(\roman*)]
\setlength{\itemsep}{0pt}
\item $X_n\xrightarrow{qm}X\Rightarrow X_n\xrightarrow{\mathbb{P}}X$.
\item $X_n\xrightarrow{\mathbb{P}}X\Rightarrow X_n\rightsquigarrow X$.
\item $X_n\rightsquigarrow c\Rightarrow X_n\xrightarrow{\mathbb{P}}c$.
\item $X_n\xrightarrow{\mathrm{a.s.}}X\Rightarrow X_n\xrightarrow{\mathbb{P}}X$.
\item $X_n\xrightarrow{L^1}X\Rightarrow X_n\xrightarrow{\mathbb{P}}X$, and $X_n\xrightarrow{qm}X\Rightarrow X_n\xrightarrow{L^1}X$.
\end{enumerate}
\begin{proof}
\begin{enumerate}[label=(\roman*)]
\setlength{\itemsep}{0pt}
\item If $X_n\xrightarrow{qm}X$, then for each $\varepsilon>0$, Markov gives $\mathbb{P}(\abs{X_n-X}>\varepsilon)\leq\varepsilon^{-2}\mathbb{E}\abs{X_n-X}^2\to0$.
\item Let $x$ be a continuity point of $F$.\\
For each $\varepsilon>0$, $F(x-\varepsilon)-\mathbb{P}(\abs{X_n-X}>\varepsilon)\leq F_n(x)\leq F(x+\varepsilon)+\mathbb{P}(\abs{X_n-X}>\varepsilon)$,\\
so $F(x-\varepsilon)\leq\liminf_nF_n(x)\leq\limsup_nF_n(x)\leq F(x+\varepsilon)$; letting $\varepsilon\downarrow0$ gives $F_n(x)\to F(x)$.
\item For $\varepsilon>0$, the distribution function of $c$ is continuous at $c-\varepsilon$ and $c+\varepsilon$.\\
Hence $\mathbb{P}(\abs{X_n-c}>\varepsilon)\leq F_n(c-\varepsilon)+1-F_n(c+\varepsilon)\to0$.
\item Fix $\varepsilon>0$ and set $Y_n=\mathbf{1}_{\{\abs{X_n-X}>\varepsilon\}}$. If $X_n\to X$ a.s., then $Y_n\to0$ a.s. and $0\leq Y_n\leq1$.\\
Hence dominated convergence gives $\mathbb{P}(\abs{X_n-X}>\varepsilon)=\mathbb{E}Y_n\to0$.
\item If $X_n\xrightarrow{L^1}X$, then $\mathbb{P}(\abs{X_n-X}>\varepsilon)\leq\varepsilon^{-1}\mathbb{E}\abs{X_n-X}\to0$ by Markov.\\
If $X_n\xrightarrow{qm}X$, then $\mathbb{E}\abs{X_n-X}\leq(\mathbb{E}\abs{X_n-X}^2)^{1/2}\to0$ by Cauchy--Schwarz.
\end{enumerate}
\hlt{Structural Commentary}: Strong modes give tail bounds; distribution convergence gives c.d.f. control.\\
\hlt{Key Insight}: Convert to tail bounds; use c.d.f. squeezes when laws enter.
\end{proof}
\end{theorem}

\begin{figure}[H]
  \centering
  % Source: supplied Chapter 6, Figure 6.1; simplified implication map.
\begin{tikzpicture}[
  >=Stealth,
  line cap=round,
  line join=round,
  every node/.style={font=\scriptsize},
  mode/.style={
    draw=blue!60!black,
    fill=blue!6,
    rounded corners=2pt,
    text=black,
    line width=0.55pt,
    minimum width=2.10cm,
    minimum height=0.66cm,
    align=center,
    inner xsep=5pt,
    inner ysep=3pt
  },
  imp/.style={->, black!78, line width=0.8pt},
  cond/.style={->, blue!70!black, dotted, line width=0.9pt}
]
  \node[mode] (qm) at (0,0) {$X_n\xrightarrow{qm}X$\\quadratic mean};
  \node[mode] (prob) at (2.95,0) {$X_n\xrightarrow{\mathbb{P}}X$\\probability};
  \node[mode] (dist) at (5.90,0) {$X_n\rightsquigarrow X$\\distribution};

  \draw[imp] (qm) -- (prob);
  \draw[imp] (prob) -- (dist);
  \draw[cond] ([yshift=-0.18cm]dist.west) -- ([yshift=-0.18cm]prob.east);
\end{tikzpicture}

  \caption{Solid: always; \textcolor{blue!70!black}{dotted}: point-mass limit.}
  \label{fig:probabilistic-convergence-modes}
\end{figure}

\begin{example} \label{ex:convergence-implication-nonconverses} \hlt{Why the Reverse Arrows Fail}
\begin{enumerate}[label=(\roman*)]
\setlength{\itemsep}{0pt}
\item Let $U\sim\operatorname{Uniform}(0,1)$ and $X_n=\sqrt{n}\mathbf{1}_{(0,1/n)}(U)$. Then for each $\varepsilon>0$, $\mathbb{P}(\abs{X_n}>\varepsilon)\leq1/n\to0$,\\
so $X_n\xrightarrow{\mathbb{P}}0$, but $\mathbb{E}X_n^2=1$ for all $n$; hence $X_n\nrightarrow0$ in quadratic mean.
\item Let $X\sim N(0,1)$ and set $X_n=-X$. Then $X_n\sim X$ for all $n$, so $X_n\rightsquigarrow X$,\\
but for every $\varepsilon>0$, $\mathbb{P}(\abs{X_n-X}>\varepsilon)=\mathbb{P}(2\abs{X}>\varepsilon)>0$ for all $n$; hence $X_n\nrightarrow X$ in probability.
\end{enumerate}
\end{example}

\begin{remark} \label{rem:convergence-in-probability-does-not-control-expectations} \hlt{A Warning About Expectations}\\
Convergence in probability does not imply convergence of expectations.\\
If $\mathbb{P}(X_n=n^2)=1/n$, $\mathbb{P}(X_n=0)=1-1/n$, then $\forall \varepsilon>0$, $\mathbb{P}(\abs{X_n}>\varepsilon)=1/n\to0$, so $X_n\xrightarrow{\mathbb{P}}0$, $\mathbb{E}X_n=n\to+\infty$.
\end{remark}

\begin{definition} \label{def:asymptotic-uniform-integrability} \hlt{Asymptotic Uniform Integrability}\\
A sequence $(X_n)$ is \hlt{asymptotically uniformly integrable} if $\lim_{M\to\infty}\limsup_{n\to\infty}\mathbb{E}(\abs{X_n}\mathbf{1}_{\{\abs{X_n}>M\}})=0$.
\end{definition}

\begin{proposition} \label{prop:aui-plus-probability-gives-expectation-convergence} \hlt{Recovering Expectations from Probability Convergence}\\
Let $b\in\mathbb{R}$. If $X_n\xrightarrow{\mathbb{P}}b$ and $(X_n)$ is asymptotically uniformly integrable, then $\mathbb{E}X_n\to b$.
\begin{proof}
For $M>|b|$, set $\tau_M(x)=(-M)\vee(x\wedge M)$. As $\tau_M$ bounded, continuous; $\tau_M(b)=b$, $\tau_M(X_n)\xrightarrow{\mathbb{P}}b$.\\
For $\eta>0$, $\mathbb{E}\abs{\tau_M(X_n)-b}\leq\eta+2M\mathbb{P}(\abs{\tau_M(X_n)-b}>\eta)$, so $\mathbb{E}\tau_M(X_n)\to b$ after $n\to\infty$ and $\eta\downarrow0$.\\
Moreover, $\abs{X_n-\tau_M(X_n)}\leq\abs{X_n}\mathbf{1}_{\{\abs{X_n}>M\}}$, hence $\limsup_n\abs{\mathbb{E}X_n-b}\leq\limsup_n\mathbb{E}(\abs{X_n}\mathbf{1}_{\{\abs{X_n}>M\}})$.\\
Letting $M\to\infty$ and using asymptotic uniform integrability gives $\mathbb{E}X_n\to b$.\\
\hlt{Structural Commentary}: Truncation makes convergence bounded; asymptotic uniform integrability kills tails.\\
\hlt{Key Insight}: Probability convergence handles the truncations; tail control restores expectations.
\end{proof}
\end{proposition}

\begin{theorem} \label{thm:continuous-mapping-and-slutsky} \hlt{Continuous Mapping and Slutsky Rules}\\
Let $g:\mathbb{R}\to\mathbb{R}$ be continuous, and let $c\in\mathbb{R}$.
\begin{enumerate}[label=(\roman*)]
\setlength{\itemsep}{0pt}
\item $X_n\xrightarrow{\mathbb{P}}X\Rightarrow g(X_n)\xrightarrow{\mathbb{P}}g(X)$.
\item $X_n\rightsquigarrow X\Rightarrow g(X_n)\rightsquigarrow g(X)$.
\item $X_n\rightsquigarrow X$ and $Y_n\xrightarrow{\mathbb{P}}c\Rightarrow X_n+Y_n\rightsquigarrow X+c$ and $X_nY_n\rightsquigarrow cX$.
\item $X_n\xrightarrow{\mathbb{P}}X$ and $Y_n\xrightarrow{\mathbb{P}}Y\Rightarrow X_n+Y_n\xrightarrow{\mathbb{P}}X+Y$ and $X_nY_n\xrightarrow{\mathbb{P}}XY$.
\end{enumerate}
\begin{proof}
\begin{enumerate}[label=(\roman*)]
\setlength{\itemsep}{0pt}
\item Fix $\varepsilon,\eta>0$. Choose $K$, $\mathbb{P}(\abs{X}>K)<\eta$, $\delta\leq1$ s.t. $\abs{x-y}<\delta$, $\abs{x},\abs{y}\leq K+1$; $\abs{g(x)-g(y)}<\varepsilon$.\\
Then $\mathbb{P}(\abs{g(X_n)-g(X)}>\varepsilon)\leq\mathbb{P}(\abs{X}>K)+\mathbb{P}(\abs{X_n-X}>\delta)$, so $\limsup$ at most $\eta$.
\item Let $t$ be a continuity point of the law of $g(X)$ and set $A_t=\{x:g(x)\leq t\}$.\\
Since $\partial A_t\subseteq\{x:g(x)=t\}$ and $\mathbb{P}(g(X)=t)=0$, the continuity-set form of weak convergence gives\\
$\mathbb{P}(g(X_n)\leq t)=\mathbb{P}(X_n\in A_t)\to\mathbb{P}(X\in A_t)=\mathbb{P}(g(X)\leq t)$.
\item If $Z_n\rightsquigarrow Z$ and $R_n\xrightarrow{\mathbb{P}}0$, then $Z_n+R_n\rightsquigarrow Z$ by the c.d.f. squeeze\\
$F_{Z_n}(t-\eta)-\mathbb{P}(\abs{R_n}>\eta)\leq F_{Z_n+R_n}(t)\leq F_{Z_n}(t+\eta)+\mathbb{P}(\abs{R_n}>\eta)$,\\
followed by $n\to\infty$ and $\eta\downarrow0$ through continuity points of $F_Z$.\\
Since $X_n+c\rightsquigarrow X+c$ by (ii), applying this with $Z_n=X_n+c$ and $R_n=Y_n-c$ gives $X_n+Y_n\rightsquigarrow X+c$.\\
Also $X_n\rightsquigarrow X$ implies $(X_n)$ is tight, so $\mathbb{P}(\abs{X_n(Y_n-c)}>\varepsilon)\leq\mathbb{P}(\abs{X_n}>K)+\mathbb{P}(\abs{Y_n-c}>\varepsilon/K)$.\\
Letting $n\to\infty$, then $K\uparrow\infty$ gives $X_n(Y_n-c)\xrightarrow{\mathbb{P}}0$.\\
Since $cX_n\rightsquigarrow cX$ by (ii), the additive perturbation just proved gives $X_nY_n=cX_n+X_n(Y_n-c)\rightsquigarrow cX$.
\item For sums, $\mathbb{P}(\abs{(X_n+Y_n)-(X+Y)}>\varepsilon)\leq\mathbb{P}(\abs{X_n-X}>\varepsilon/2)+\mathbb{P}(\abs{Y_n-Y}>\varepsilon/2)\to0$.\\
For products, $(X_n)$ and $Y$ are tight, $\mathbb{P}(\abs{X_n(Y_n-Y)}>\varepsilon/2)\leq\mathbb{P}(\abs{X_n}>K)+\mathbb{P}(\abs{Y_n-Y}>\varepsilon/(2K))$,\\
while $\mathbb{P}(\abs{Y(X_n-X)}>\varepsilon/2)\leq\mathbb{P}(\abs{Y}>K)+\mathbb{P}(\abs{X_n-X}>\varepsilon/(2K))$; let $n\to\infty$ and then $K\uparrow\infty$.
\end{enumerate}
\hlt{Structural Commentary}: Continuity transports convergence; Slutsky absorbs probability-small errors.\\
\hlt{Key Insight}: Control the bounded high-probability part; let the tails vanish.
\end{proof}
\end{theorem}
